% Options for packages loaded elsewhere
\PassOptionsToPackage{unicode}{hyperref}
\PassOptionsToPackage{hyphens}{url}
%
\documentclass[10pt]{article}
\usepackage{times}
\usepackage[top=0.3in,left=0.3in,right=0.3in,bottom=0.4in,landscape,a4paper,
footskip=0.2in]{geometry}
% \usepackage[handout]{fds}
% \todosfalse
\usepackage[math]{iwona}
\usepackage{hyperref}

\usepackage{amsmath,amssymb}
\usepackage{iftex}
\ifPDFTeX
  \usepackage[T1]{fontenc}
  \usepackage[utf8]{inputenc}
  \usepackage{textcomp} % provide euro and other symbols
\else % if luatex or xetex
  \usepackage{unicode-math} % this also loads fontspec
  \defaultfontfeatures{Scale=MatchLowercase}
  \defaultfontfeatures[\rmfamily]{Ligatures=TeX,Scale=1}
\fi
% Use upquote if available, for straight quotes in verbatim environments
\IfFileExists{upquote.sty}{\usepackage{upquote}}{}
\usepackage{xcolor}
\usepackage{longtable,booktabs,array}
\usepackage{calc} % for calculating minipage widths
% Correct order of tables after \paragraph or \subparagraph
\usepackage{etoolbox}
\makeatletter
\patchcmd\longtable{\par}{\if@noskipsec\mbox{}\fi\par}{}{}
\makeatother
% Allow footnotes in longtable head/foot
\IfFileExists{footnotehyper.sty}{\usepackage{footnotehyper}}{\usepackage{footnote}}
\makesavenoteenv{longtable}
\setlength{\emergencystretch}{3em} % prevent overfull lines
\providecommand{\tightlist}{%
  \setlength{\itemsep}{0pt}\setlength{\parskip}{0pt}}
\setcounter{secnumdepth}{-\maxdimen} % remove section numbering
\ifLuaTeX
  \usepackage{selnolig}  % disable illegal ligatures
\fi

\usepackage{enumitem}
\setlist{topsep=0pt,leftmargin=10pt,noitemsep,parsep=0pt}
\usepackage{fancyhdr}

% For callout boxes 
\usepackage[skins,breakable]{tcolorbox}
% \@ifpackageloaded{tcolorbox}{}{\usepackage[skins,breakable]{tcolorbox}}
\usepackage{fontawesome5}

% \newcommand{\principle}[2]{\paragraph{#1} #2 \\[0.2em]\hrule \\[0.2em]}

\definecolor{quarto-callout-tip-color}{HTML}{00A047}
\definecolor{quarto-callout-tip-color-frame}{HTML}{02b875}
\definecolor{quarto-callout-standard-color}{HTML}{909090}
\definecolor{quarto-callout-standard-color-frame}{HTML}{acacac}

\author{}
\date{}

\begin{document}
\fancyfoot[C]{
  \footnotesize Version 2.0-beta2 \copyright 2022-2025 by David Sterratt, licensed under CC
  BY-SA 4.0
  Source: \url{https://github.com/Inf2-FDS/fds-visualisation}}
\thispagestyle{fancy}

\begin{tcolorbox}[leftrule=.15mm,rightrule=.15mm,
  toprule=.15mm,bottomrule=.15mm,arc=.35mm,
  colback=quarto-callout-standard-color!35!white,
  enhanced jigsaw,
  drop shadow,
  colframe=quarto-callout-standard-color-frame,top=1mm,bottom=1mm,
  halign=flush center]
  \large\textbf{Informatics 2 -- Foundations of Data Science: Criteria for marking
visualisation explanations/interpretations and
code}
\end{tcolorbox}

\begin{longtable}[]{@{}
  >{\raggedright\arraybackslash\bfseries}p{(\columnwidth - 4\tabcolsep) * \real{0.1366}}
  >{\raggedright\arraybackslash}p{(\columnwidth - 4\tabcolsep) * \real{0.4899}}
  >{\raggedright\arraybackslash}p{(\columnwidth - 4\tabcolsep) * \real{0.3735}}@{}}
\toprule\noalign{}
\begin{minipage}[b]{\linewidth}\raggedright
\textbf{Mark}
\end{minipage} & \begin{minipage}[b]{\linewidth}\raggedright
\textbf{Explanation/interpretation}
\end{minipage} & \begin{minipage}[b]{\linewidth}\raggedright
\textbf{Code readability}
\end{minipage} \\
\midrule\noalign{}
\endhead
\bottomrule\noalign{}
\endlastfoot
4 - Excellent & \begin{minipage}[t]{\linewidth}\raggedright
\begin{itemize}
\item
  Easy to comprehend due to clear structure and precision.
\item
  Clear and correct summarisation of multiple patterns and noteworthy
  features, possibly with sensible explanations or hypotheses for them
  or relating them to other data or knowledge.
\item
  Interpretation means the visualisation makes sense.
\end{itemize}
\end{minipage}

& \begin{minipage}[t]{\linewidth}\raggedright
Code is very easy to read, due to

\begin{itemize}
\item
  code being clean with appropriate comments
\item
  appropriate use of pandas functions (such as pd.merge() and
  .groupby()).
\end{itemize}
\end{minipage} \\
\hline
3 - Good & \begin{minipage}[t]{\linewidth}\raggedright
\begin{itemize}
\item
  Straightforward to comprehend due to reasonably clear structure and
  precision.
\item
  Clear and correct summarisation of some patterns and noteworthy
  features.
\item
  Interpretation has added significantly to the visualisation.
\end{itemize}
\end{minipage} & \begin{minipage}[t]{\linewidth}\raggedright
Code is easy to read

\begin{itemize}
\item
  due to code being clean
\item
  appropriate use of pandas functions (such as pd.merge() and
  .groupby()).
\end{itemize}
\end{minipage}

\\
\hline
2 - Fair & \begin{minipage}[t]{\linewidth}\raggedright
\begin{itemize}
\item
  Somewhat difficult to comprehend due to structure or vagueness
\item
  but is still clear that a few patterns and noteworthy features have
  been summarised or identified correctly, though there may be minor
  errors.
\item
  Interpretation has added a bit to the visualisation.
\end{itemize}
\end{minipage} & \begin{minipage}[t]{\linewidth}\raggedright
Code is understandable with some work due to EITHER:

\begin{itemize}
\item
  Not being clean, for example with commented out or unnecessary code
\item
  but still using pandas functions appropriately
\end{itemize}

OR

\begin{itemize}
\item
  Being clean
\item
  but using multiple lines of code to achieve what could have been done
  with pandas functions.
\end{itemize}
\end{minipage}

\\
\hline
1 - Inadequate & \begin{minipage}[t]{\linewidth}\raggedright
\begin{itemize}
\item
  Very difficult to comprehend, due to poor structure or vagueness,
\item
  but at least one pattern appears to have been summarised correctly or
  a noteworthy feature identified.
\item
  Maybe some significant errors.
\item
  Interpretation has not added to visualisation.
\end{itemize}
\end{minipage}

& \begin{minipage}[t]{\linewidth}\raggedright
Code is difficult to understand due to

\begin{itemize}
\item
  not being clean and unclear code
\item
  probably not using pandas functions.
\end{itemize}
\end{minipage} \\
\hline
0 - Absent & No meaningful interpretation. & Code is clearly not
functional.

\\
\end{longtable}

\end{document}

% LocalWords:  unicode hyperref url groupby
